% ============================================================
% XeLaTeX + acmart + 中文 兼容性修复 (必须在\documentclass之前)
% ============================================================
% 1. 阻止 fontenc 加载 (与XeLaTeX的fontspec冲突)
\makeatletter
\@namedef{ver@fontenc.sty}{}
\makeatother

% 2. 禁用microtype (与XeLaTeX+ctex冲突)
\PassOptionsToPackage{activate=false}{microtype}

% 3. 禁用inconsolata和 libertine (XeLaTeX用系统字体)
\PassOptionsToPackage{type1}{inconsolata}
% ============================================================

\documentclass[sigconf]{acmart}

% ---- 中文支持 ----
% scheme=plain: 仅提供中文支持, 不改变文档类格式 (避免与acmart冲突)
\usepackage[scheme=plain]{ctex}

% ---- Packages ----
% acmart已加载: amsmath, graphicx, hyperref
% amsmath, amssymb, amsfonts 也已由acmart加载
\usepackage{booktabs}
\usepackage{multirow}
\usepackage{algorithm}
\usepackage{algpseudocode}

% ---- 中文适配 ----
\renewcommand{\abstractname}{摘要}
\renewcommand{\keywordsname}{关键词}
\renewcommand{\figurename}{图}
\renewcommand{\tablename}{表}
\renewcommand{\refname}{参考文献}
\floatname{algorithm}{算法}
\renewcommand{\algorithmicrequire}{\textbf{输入:}}
\renewcommand{\algorithmicensure}{\textbf{输出:}}
\renewcommand{\listalgorithmname}{算法列表}

% ---- cleveref (最后加载) ----
\usepackage{cleveref}
\crefname{equation}{式}{式}
\crefname{figure}{图}{图}
\crefname{table}{表}{表}
\crefname{section}{节}{节}
\crefname{algorithm}{算法}{算法}

% Math macros
% 共享数学宏

% 图符号
\newcommand{\graph}{\mathcal{G}}
\newcommand{\chain}{\mathcal{C}}
\newcommand{\event}{e}
\newcommand{\eventset}{\mathcal{E}}

% 度量符号
\newcommand{\phiMetric}{\phi}
\newcommand{\chainCoherence}{\Phi}
\newcommand{\branchCoherence}{\Phi_{\text{bw}}}
\newcommand{\hoplimit}{\delta}
\newcommand{\timedecay}{\beta}
\newcommand{\graphdecay}{\alpha}
\newcommand{\branchpenalty}{\lambda}
\newcommand{\timegap}{\Delta t}
\newcommand{\graphdist}{d}

% 统计符号
\newcommand{\mean}{\mu}
\newcommand{\std}{\sigma}

% 数学运算符
\DeclareMathOperator*{\argmax}{arg\,max}
\DeclareMathOperator*{\argmin}{arg\,min}
\DeclareMathOperator{\cosine}{cos}
\DeclareMathOperator{\branchfactor}{BF}


% ---- 文档信息 ----
\title{基于因果边感知的时序链提取方法用于溯源图APT检测}

\author{匿名作者}
\affiliation{匿名机构}

\begin{document}

\maketitle

% §0 摘要

高级持续性威胁（APT）以多步因果链的形式展开——初始入侵导致持久化，持久化导致横向移动，最终导致数据窃取。然而，基于溯源图的入侵检测系统虽然在显式编码因果依赖关系的图上运行，却仍然以粗粒度的时间窗口、孤立实体或图快照来检测攻击。没有任何方法能够识别构成攻击叙事的实际因果事件序列。我们提出了首个以因果事件链为粒度的APT检测系统。我们的核心贡献是一个模型无关的因果相干性度量（$\branchCoherence$），它量化从溯源图中提取的事件链的因果结构完整性，其理论基础是操作系统的核心原理：进程是独占的能力（capability），而文件和套接字是共享资源。基于肘部检测的校准协议仅从良性数据中自动发现操作系统的自然因果视野，无需攻击标签，也无需人工参数调优。在$\branchCoherence$指导下提取的因果链，由一个仅在良性事件序列上以自回归方式训练的自监督Transformer进行验证。在DARPA透明计算基准上，我们的因果链达到了[X]\%的链级F1分数，超过基于窗口（EagleEye）、基于随机游走（Sentient）和基于时序子图（GET-AID）的替代方法超过10个百分点的绝对F1提升，同时相比最先进的窗口级检测器（KAIROS）减少了超过50倍的调查成本。校准协议在不同的操作系统配置上发现了不同的因果视野，展示了跨平台自动适配能力。

% §1 引言

高级持续性威胁（APT）不是孤立的事件——它们是因果链。攻击者利用Web服务器的漏洞，派生出Shell，下载载荷，读取凭证，将数据外泄到外部主机。每一步\emph{导致}下一步。对于调查告警的安全分析员来说，知道\emph{哪条因果相连的事件序列}构成了攻击，比知道哪个15分钟时间窗口或哪个单独文件可疑更具可操作性。

溯源图由内核级审计日志构建，天然捕获了这些因果依赖关系。节点代表系统实体（进程、文件、套接字），有向边代表带有精确时间戳的因果系统事件（进程派生、文件写入、网络发送）。由于内核调解所有系统交互，溯源图提供了完整且防篡改的因果系统活动记录。

尽管运行在显式编码因果的数据结构上，基于溯源的入侵检测系统（PIDS）尚未实现因果链检测的潜力：

\begin{enumerate}
\item \textbf{粗粒度检测。} 最先进的PIDS在时间窗口~\cite{cheng2024kairos}、单个实体~\cite{jia2024magic,wang2022threatrace}、图快照~\cite{yang2023prographer}或边级别~\cite{edgetrace2025}进行检测。没有任何方法识别哪条具体的因果相连的事件序列构成了攻击。
\item \textbf{启发式链提取。} 最近的Transformer-on-provenance方法使用固定时间窗口~\cite{eagleeye2024}、随机游走~\cite{sentient2026}或时序子图~\cite{getaid2025}来分割溯源图。它们都不尊重因果边方向——固定窗口混合了因果无关的事件，随机游走可能沿逆向遍历边，时序子图使用时间而非因果作为组织原则。
\item \textbf{后处理重建。} 产生攻击路径的方法是在节点/边级检测之后作为后处理步骤进行的~\cite{slot2025,cage2025}。链质量从根本上受限于上游检测的召回率——漏掉一个关键节点就会导致链断裂。
\end{enumerate}

一个更深层的问题潜藏于所有这些局限之下：\textbf{没有任何已有方法在提取链之前定义了什么是“因果上好”的链。}已有方法将链提取视为一个分割问题——如何将大图切成小片段——并依赖下游模型性能间接验证分割质量。

\textbf{我们的方法。} 我们将链提取重新定义为\emph{度量}问题。我们提出一个模型无关的度量，直接量化从溯源图中提取的任意事件链的因果结构完整性，以及一个仅从良性数据中发现该度量参数的校准协议。

具体而言，我们做出了以下贡献：

\begin{enumerate}
\item \textbf{因果相干性度量（$\branchCoherence$）。} 一个模型无关的度量，根植于操作系统原理：进程是独占的能力（内核保证进程身份），而文件和套接字是共享资源（受第三方修改影响）。该度量采用三分法的逐对相干性函数和惩罚高分支良性进程的结构稀疏先验。
\item \textbf{无监督校准协议。} 基于肘部检测的方法，仅从良性溯源数据中发现操作系统的自然因果视野，零攻击标签、零人工参数调优。该协议具有跨平台性：自动适配不同操作系统之间审计日志粒度的差异。
\item \textbf{端到端检测系统。} 流水线组合了基于TGN的种子识别（复用KAIROS）、以$\branchCoherence$为指导的因果链提取器、以及通过自回归预测事件过渡来区分攻击链与良性链的自监督Transformer。Transformer仅在良性数据上训练。
\item \textbf{实证验证。} 在DARPA透明计算基准上，因果链达到[X]\%的链级F1，超过基于窗口、基于随机游走、基于子图和基于规则的链提取方法超过10个百分点的绝对F1，同时相比KAIROS减少超过50倍的调查成本。校准协议在不同数据集（E3和E5）上发现不同的因果视野，展示了跨平台自动适配能力。
\end{enumerate}

本文其余部分组织如下。\S\ref{sec:background}提供溯源图背景并分析已有链提取策略的局限性。\S\ref{sec:related}将我们的工作置于更广泛的溯源入侵检测文献中。\S\ref{sec:design}展示系统设计，详细讨论因果链提取器、相干性度量和自回归Transformer。\S\ref{sec:metric}在独立于任何模型训练的情况下验证相干性度量。\S\ref{sec:evaluation}呈现端到端检测结果。\S\ref{sec:discussion}讨论部署考量和局限性。\S\ref{sec:conclusion}总结。

% §2 背景与动机
\label{sec:background}

\subsection{溯源图与APT检测}

溯源图 $\graph = (V, E)$ 是由内核级审计日志构建的有向图。每个节点 $v \in V$ 代表一个系统实体（进程、文件或网络套接字），每条有向边 $e = (u \rightarrow v, t, \tau) \in E$ 代表一个因果系统事件——实体 $u$ 在时间戳 $t$ 通过事件类型 $\tau \in \{\text{FORK}, \text{EXEC}, \text{READ}, \text{WRITE}, \text{SENDTO}, \text{RECVFROM}, \dots\}$ 影响了实体 $v$。

溯源图具备两个使其在入侵检测中格外强大的特性。第一，它们是\emph{完备的}：由于内核调解所有系统交互，系统实体之间的每个因果关系都被记录下来。第二，它们是\emph{防篡改的}：攻击者无法在未获得内核级访问权限的情况下修改内核级审计日志——而一旦获得内核级访问，系统已经彻底沦陷。

溯源数据的规模带来了根本性挑战。DARPA透明计算E3数据集（一个标准基准）包含268,242个节点和29,727,441条边，跨越五天。从这样的体量中提取有意义的攻击信号需要有原则的过滤。

\subsection{检测粒度鸿沟}

已有PIDS在逐步细化的粒度上运行，形成一个谱系：

\begin{itemize}
\item \textbf{图级} (Unicorn~\cite{unicorn2020})：将整个溯源图分类为良性或恶意。
\item \textbf{快照级} (PROGRAPHER~\cite{yang2023prographer})：将图划分为时序快照并对每个快照进行分类。
\item \textbf{窗口级} (KAIROS~\cite{cheng2024kairos})：在15分钟时间窗口内对事件评分，标记异常密度高的窗口。
\item \textbf{实体级} (MAGIC~\cite{jia2024magic}, ThreaTrace~\cite{wang2022threatrace})：对单个系统实体（进程、文件）评分。
\item \textbf{边级} (EdgeTrace~\cite{edgetrace2025})：对单条溯源边进行异常评分。
\end{itemize}

每个级别回答了一个逐步细化的问题，但没有人回答对调查最重要的那个问题：\textbf{“哪条因果相连的事件序列构成了攻击？”} 我们称之为\emph{链级}检测。

\subsection{已有链提取方法的局限性}

近期的工作已经开始将Transformer应用于溯源图。尽管这些方法代表了重要进展，但它们的链提取策略共享一个共同的局限性：不尊重因果边方向。

\textbf{固定窗口 (EagleEye~\cite{eagleeye2024})。} EagleEye将溯源图分割为固定大小的时间窗口，并将每个窗口内的事件线性化为Transformer分类的序列。该方法有两个缺陷。首先，时间$T$处的固定窗口可能包含恰好同时发生的、来自不相关进程的事件——Firefox页面加载和SSH登录可能落入同一窗口但毫无因果关系。其次，跨越窗口边界的因果相连事件被切断。

\textbf{随机游走 (Sentient~\cite{sentient2026})。} Sentient使用基于随机游走的场景分割来为Graph Transformer提取子图。随机游走在本质上是方向无关的：一次游走可能沿逆向遍历溯源边，将读取文件的过程与写入文件的过程连接起来，而不尊重这些事件发生的时间顺序。

\textbf{时序子图 (GET-AID~\cite{getaid2025})。} GET-AID基于事件时间戳的接近程度构造时序子图。虽然时间戳提供了一个弱信号，但时间接近并不意味因果连接：两个同时发生的事件可能属于完全独立的进程树。

\textbf{后处理重建 (SLOT~\cite{slot2025}, CAGE~\cite{cage2025})。} 这些方法首先检测异常节点或边，然后通过聚类或因果加权将它们连接成攻击路径。链是后处理的产出物，而非检测单元本身。关键的是，如果上游检测器漏掉单个关键节点，重建的链就不完整——且没有机制可以恢复它。

\textbf{基于规则的提取 (SPARSE~\cite{sparse2024}, Holmes~\cite{holmes2019})。} 这些方法使用专家定义的规则和威胁情报对路径进行评分和提取。它们无法检测其模式不在预定义规则集内的新型攻击。

\subsection{缺失的机制}

所有上述方法共享一个基本做法：首先决定\emph{如何分割}溯源图（按时间、按随机游走、按规则），然后依赖下游模型来验证分割是否有用。没有方法首先定义\emph{什么属性使一条链因果相干}，然后提取满足这些属性的链。

这是我们要填补的缺口。在\S\ref{sec:design}中，我们定义了一个量化链质量的因果相干性度量（独立于任何模型），并设计了一个最大化该度量的提取算法。

% 相关工作
\section{相关工作}
\label{sec:related}

我们按 APT 检测逻辑组织相关工作，而非按具体模型架构分类。现有研究大体可以分为两类：一类从异常行为出发，学习正常系统行为并寻找偏离；另一类从攻击过程出发，试图恢复攻击路径或行为序列。本文工作与两类方法都相关：我们借鉴异常检测方法发现可疑起点，但把最终检测对象提升为因果链，并用链级模型判断该过程是否异常。

\paragraph{异常行为驱动的 APT 检测。}

第一类工作将 APT 检测视为异常发现问题：先学习良性系统行为的表示，再在图、时间窗口、节点、实体或边的粒度上寻找偏离。UNICORN~\cite{unicorn2020} 使用 Weisfeiler--Lehman 图核和直方图草图对整个溯源图进行分类；PROGRAPHER~\cite{yang2023prographer} 将溯源图划分为时序快照，用 graph2vec 和 TextRCNN 进行序列级异常检测。MAGIC~\cite{jia2024magic} 通过图掩码自编码器学习良性实体表示，并用 KNN 密度估计发现异常实体；ThreaTrace~\cite{wang2022threatrace} 和 FLASH~\cite{flash2024} 分别使用 GraphSAGE 或语义/图上下文编码进行节点级检测。KAIROS~\cite{cheng2024kairos} 进一步在边级计算时序图网络重建损失，并在时间窗口内聚合异常信号。EdgeTrace~\cite{edgetrace2025} 也以边级异常评分作为后续追踪的基础。

这类方法的优势是能够在大规模审计数据中有效定位“哪里不正常”：异常时间段、可疑实体、可疑节点或可疑边。但其检测对象通常仍是局部或聚合单元，而不是攻击调查真正需要的完整行为过程。因此，分析员往往还需要在告警之后手工判断这些局部异常之间是否存在因果联系。本文继承这一路线的异常检测能力：在当前实现中，我们使用 KAIROS 的边级评分作为种子来源；但不同之处在于，边级异常只用于提示“从哪里开始查”，而最终检测和评估发生在由因果约束提取出的事件链上。

\paragraph{攻击过程驱动的检测与调查。}

第二类工作直接面向攻击过程，试图把分散事件组织成路径、序列或子图。Holmes~\cite{holmes2019}、Poirot~\cite{poirot2019} 和 SLEUTH 通过专家定义的攻击模式在溯源图中追踪因果路径；SPARSE~\cite{sparse2024} 结合语义标签和威胁情报构造可疑语义图。近期学习方法也开始关注路径或序列层面的检测：SLOT~\cite{slot2025} 在异常检测后用标签传播聚类构造攻击链；CAGE~\cite{cage2025} 通过 Q-learning 为溯源边赋予动态因果权重，并提取加权攻击子图；EdgeTrace~\cite{edgetrace2025} 在边级异常评分之后进行语义聚类和路径推断。另一条相关路线是把溯源数据组织成序列后交给 Transformer 或类似架构建模，例如 EagleEye~\cite{eagleeye2024} 的固定窗口事件序列、Sentient~\cite{sentient2026} 的图 Transformer 与意图分析模块、GET-AID~\cite{getaid2025} 的时序子图注意力，以及 PanThreat~\cite{panthreat2025} 和 LogShield~\cite{logshield2024} 的日志/溯源序列建模。

这类方法更接近攻击调查需求，因为它们试图回答“攻击如何展开”。但其关键差异在于链或序列如何产生：有的方法依赖人工攻击模板或威胁情报，难以覆盖未知 APT；有的方法把路径重建放在节点/边级告警之后，链质量受上游漏报限制；还有的方法用时间窗口、随机游走或时序子图定义序列，未必保证相邻事件之间具有真实依赖关系。换言之，这些方法已经认识到过程级建模的重要性，但过程本身往往由模板、分割规则或后处理决定，而不是由一个可计算的因果链质量准则在检测前显式定义。

\paragraph{与本文工作的关系。}

本文属于攻击过程驱动的检测与调查方向，目标是把 APT 告警从局部异常推进到可解释的行为过程。与仅输出异常窗口、实体或边的方法不同，本文将最终检测对象定义为因果事件链；与依赖人工模板或告警后路径重建的方法不同，本文在检测之前提取结构可信的因果链，并在链级判断其是否偏离良性行为。本文的主要亮点在于把“因果链提取”和“链级异常检测”明确分层：前者保证检测对象具有可解释的依赖结构，后者利用自监督 Transformer 学习正常链级行为并识别语义异常。具体方法将在 \S\ref{sec:method} 中介绍。

% §4 系统设计
\label{sec:design}

\subsection{概述}

图~\ref{fig:architecture}展示了系统架构。流水线由离线阶段（每次部署运行一次）和在线阶段（每个测试日运行）组成。

在离线阶段，我们（1）从完整溯源图计算逐节点属性（Node Profiler），（2）在良性数据上预训练TGN进行边类型预测（复用KAIROS），（3）通过良性链上的肘部检测校准因果相干性度量参数，（4）仅在良性链上训练自回归Transformer。

在线阶段，对于测试日：（1）TGN顺序处理事件，维护时序记忆状态，对每条边输出重建损失；（2）损失超过阈值（良性损失95\%分位数）的边触发因果链提取；（3）链提取器执行双向因果追溯，通过$\branchCoherence$剪枝分支，输出候选链；（4）自回归Transformer计算逐事件重建损失，通过预先计算的eCDF映射为异常百分位数；（5）高异常分数的链被标记，共享节点且时间接近的高分链合并为完整的攻击路径。

\subsection{Node Profiler：全局图上下文}

我们从完整溯源图中预计算五个节点级属性，并将其与KAIROS已有的16维层次路径哈希（node2higvec）拼接，产生21维的节点表示。

\begin{enumerate}
\item \textbf{逆文档频率（IDF）。} $\text{IDF}(v) = \log(T / (t_v + 1))$，其中$T$是时间窗口数，$t_v$是包含节点$v$的窗口数。IDF越高表示节点越稀有。
\item \textbf{路径敏感性。} 二值指示器：节点路径是否包含安全敏感关键词（\texttt{/etc/passwd}、\texttt{/etc/shadow}、\texttt{authorized\_keys}、\texttt{/root/}等）。
\item \textbf{度异常。} $|\text{deg}(v) - \mu_{\text{deg}}(\text{type}(v))| / \sigma_{\text{deg}}(\text{type}(v))$，衡量节点连接模式相对于其类型的异常程度。
\item \textbf{桥接中心性。} 归一化介数中心性，捕获节点在溯源图中作为结构桥梁的角色。
\item \textbf{社区ID。} 在无向溯源图上的Louvain社区划分，提供粗粒度的功能分组。
\end{enumerate}

Node Profiler是特征工程，不声称贡献。已有的Transformer-on-provenance方法（EagleEye、Sentient、GET-AID）均未系统性地将全局图统计量注入token特征。

\subsection{TGN种子识别}

我们复用KAIROS预训练的时序图网络（TGN），包括GraphAttentionEmbedding模块（2层TransformerConv，8个注意力头）、使用IdentityMessage和LastAggregator的TGNMemory模块、以及LinkPredictor MLP。TGN在良性日（D2-D4）上训练，预测每条溯源边的类型。

推理时，TGN按顺序处理事件，每个事件后更新节点记忆状态。对于每条边$e = (u \rightarrow v, t, \tau)$，输出重建损失$\mathcal{L}_{\text{TGN}}(e)$，衡量模型预测边类型$\tau$的准确性。$\mathcal{L}_{\text{TGN}}(e)$超过良性损失分布95\%分位数的边被指定为链提取的\emph{种子}。

TGN的时序记忆至关重要：节点状态随系统运行而演化，因此模型学会区分\texttt{nginx}的fork比\texttt{bash}的fork更罕见，以及某些事件类型过渡对特定节点是异常的。静态特征无法捕获这种时序上下文。

\subsection{因果链提取器}

给定种子边$s = (u_0 \rightarrow v_0, t_0, \tau_0)$和溯源图$\graph$，链提取器执行双向因果追溯以构造候选链。

\subsubsection{提取算法}

算法~\ref{alg:extract}展示了提取过程。算法维护一个候选链集合，初始仅包含种子。它交替进行后向追溯（什么导致了种子？）和前向追溯（种子导致了什么？），每一步扩展链。

\begin{algorithm}[t]
\caption{因果链提取}
\label{alg:extract}
\begin{algorithmic}[1]
\Require 溯源图 $\graph$，种子边 $s$，参数 $\hoplimit^*, \timedecay^*, \branchpenalty^*, B=20$
\Ensure 带有 $\branchCoherence$ 分数的候选链集合
\State $\text{chains} \gets \{ [s] \}$
\State $\text{current} \gets \{ s \}$
\For{$\text{direction} \in \{\text{backward}, \text{forward}\}$}
  \For{$\text{depth} = 1, 2, \dots$}
    \State $\text{next} \gets \emptyset$
    \For{each $(u \rightarrow v) \in \text{current}$}
      \If{$\text{direction} = \text{backward}$}
        \State $\text{next} \gets \text{next} \cup \textsc{IncomingCausal}(\graph, u, \hoplimit^*)$
      \Else
        \State $\text{next} \gets \text{next} \cup \textsc{OutgoingCausal}(\graph, v, \hoplimit^*)$
      \EndIf
    \EndFor
    \If{$\text{next} = \emptyset$} \State \textbf{break} \Comment{到达因果边界}
    \EndIf
    \State $\text{chains} \gets \textsc{BranchExtend}(\text{chains}, \text{next}, \text{direction})$
    \If{$|\text{chains}| > B$}
      \State $\text{chains} \gets \textsc{PruneByCoherence}(\text{chains}, B, \branchCoherence)$
    \EndIf
    \State $\text{current} \gets \text{next}$
  \EndFor
\EndFor
\State \Return $\text{chains}$
\end{algorithmic}
\end{algorithm}

\textsc{IncomingCausal}和\textsc{OutgoingCausal}在$\hoplimit^*$跳内沿有向溯源边执行BFS，过滤为表示因果影响的事件类型（EXEC、FORK、WRITE、SENDTO），排除仅表示后果的类型（READ、CLOSE）。\textsc{BranchExtend}创建已有链与新边的笛卡尔积，保持时序顺序。\textsc{PruneByCoherence}保留$\branchCoherence$排名前$B$的链。

\subsubsection{因果相干性度量}

该度量量化事件链的因果结构完整性。它是模型无关的——可直接从溯源图拓扑和时间戳计算，无需任何训练好的模型。

\textbf{逐对相干性。} 对于链中连续事件$e_i, e_{i+1}$，令$\graphdist_i$为$\graph$中从$\text{dst}(e_i)$到$\text{src}(e_{i+1})$的最短有向路径长度（若无有向路径则$\graphdist_i = \infty$），$\timegap_i = t_{i+1} - t_i$。

逐对相干性$\phiMetric(e_i, e_{i+1})$根据\emph{桥接节点类型}——当$\graphdist_i = 0$时，同时是$e_i$终点和$e_{i+1}$起点的节点——区分三种情况：

\begin{equation}
\phiMetric(e_i, e_{i+1}) =
\begin{cases}
\exp(-\graphdecay^* \cdot \graphdist_i) & \text{若 } \graphdist_i = 0 \land \text{type}(\text{dst}(e_i)) = \text{subject} \\[4pt]
\exp(-\graphdecay^* \cdot \graphdist_i) \cdot \exp(-\timedecay^* \cdot \timegap_i) & \text{若 } \graphdist_i = 0 \land \text{type}(\text{dst}(e_i)) \in \{\text{file}, \text{netflow}\} \\[4pt]
\exp(-\graphdecay^* \cdot \graphdist_i) \cdot \exp(-\timedecay^* \cdot \timegap_i) & \text{若 } 0 < \graphdist_i \leq \hoplimit^* \\
0 & \text{若 } \graphdist_i > \hoplimit^*
\end{cases}
\label{eq:phi}
\end{equation}

这三种情况根植于操作系统原理，而非任意选择：

\begin{enumerate}
\item \textbf{身份保持桥接}（$\graphdist_i = 0$，节点为进程）。进程的PID和内存空间是内核保证的独占能力。被事件$e_i$派生的进程正是发起事件$e_{i+1}$的同一个进程。因果关系是绝对的，与经过的时间无关。我们设置$\timedecay^* \approx 0$（无时间衰减）。

\item \textbf{有状态桥接}（$\graphdist_i = 0$，节点为文件或套接字）。文件和套接字是共享资源。在$e_i$的写入和$e_{i+1}$的读取之间，第三方可能已经修改了资源，打破了因果链。时间接近性提供了因果关系的证据：间隔越长，假定的因果关系越不可靠。

\item \textbf{间接路径}（$0 < \graphdist_i \leq \hoplimit^*$）。事件通过中间节点连接。因果连接是概率性的，适用时间衰减。
\end{enumerate}

\textbf{链级相干性。} 链相干性$\chainCoherence(\chain)$是逐对分数的算术平均：

\begin{equation}
\chainCoherence(\chain) = \frac{1}{n-1} \sum_{i=1}^{n-1} \phiMetric(e_i, e_{i+1})
\label{eq:chain_phi}
\end{equation}

我们使用算术平均而非几何平均是为了鲁棒性：单个缺失的审计日志事件（在生产系统中常见）会使几何平均归零。

\textbf{分支加权相干性。} 正常操作系统进程呈现幂律分支：一个Firefox进程可能写入50个文件，而编译任务派生数百个子进程。攻击链则相反，在结构上是稀疏的——DARPA E3中的nginx漏洞利用链的分支因子约为2.5。为利用这一结构性差异，我们引入了\emph{结构稀疏先验}：

\begin{equation}
\branchCoherence(\chain) = \chainCoherence(\chain) \cdot \exp(-\branchpenalty^* \cdot \branchfactor(\chain))
\label{eq:phi_bw}
\end{equation}

其中$\branchfactor(\chain)$是链中非终端节点在完整溯源图中的平均出度。参数$\branchpenalty^*$从良性数据自动校准。

\subsubsection{校准协议}

度量参数$\hoplimit^*, \timedecay^*, \graphdecay^*, \branchpenalty^*$必须在没有攻击标签的情况下设置。我们提出一个肘部检测协议，从良性数据中发现这些参数。

\begin{enumerate}
\item 使用故意宽松的参数（$\hoplimit=10$, $\timedecay=0.01$），从$K$个良性TGN种子中提取链。
\item 对每个$\hoplimit \in \{1, 2, \dots, 10\}$：从\emph{相同}的种子重新提取链（仅变动$\hoplimit$），计算所有提取链的均值$\chainCoherence$。
\item 绘制$\chainCoherence$ vs.~$\hoplimit$曲线。曲线初始上升（更多跳捕获真实因果），然后平台化或下降（噪声被强行拉入链中）。选择肘部的$\hoplimit^*$——最大负二阶导数处。
\item 固定$\hoplimit^*$；通过$\chainCoherence$ vs.~最大时间间隔$\tau$的肘部类似地校准$\timedecay^*$。
\item 设置$\graphdecay^* = 1/\hoplimit^*$, $\timedecay^* = 1/\tau^*$。
\item 将$\branchpenalty^*$校准为使得前10\%最相干的良性链的均值$\branchfactor$低于指定分位数的值。
\end{enumerate}

肘部的存在是有保证的，因为操作系统具有有限的因果深度。随着$\hoplimit$增加，提取最终将所有链连接到\texttt{init}（PID 1）或核心内核线程，将无关事件拉入导致相干性下降。

该协议是核心贡献：它零攻击标签、零人工调优，并自动适配不同OS配置（Linux auditd vs.~Windows ETW产生的日志粒度不同，因此自然因果视野也不同）。

\subsection{自回归Transformer}

提取的链由一个自监督的自回归Transformer进行验证。我们强调Transformer不是贡献——它作为链质量的测量工具。

\subsubsection{为何自回归}

我们采用自回归的next-event prediction，而非掩码事件建模（MLM/BERT）。在溯源图中，因果严格随时间向前流动。MLM允许模型在重建被掩码的过去事件时以前向的未来事件为条件——违背了因果箭头。自回归预测——每个事件仅以其因果前驱为条件——保留了溯源数据的因果语义。

\subsubsection{输入表示}

每个事件$e_i = (u \rightarrow v, t, \tau)$被编码为75维特征向量：

\begin{itemize}
\item \textbf{源节点特征}（24维）：node2higvec (16)，节点类型one-hot (3)，IDF (1)，路径敏感性 (1)，度异常 (1)，社区ID (1)，桥接中心性 (1)。
\item \textbf{边特征}（20维）：边类型one-hot (7)，$\timegap_i$正弦时间编码 (8)，$\graphdist_i$的图距离学习嵌入 (4)，桥接类型标志 (3)，是否为种子 (1)。
\item \textbf{目标节点特征}（24维）：与源节点相同结构。
\item \textbf{结构标记}（3维）：是否为分支点，是否为汇聚点，距种子深度。
\end{itemize}

一个可学习的\texttt{[BOS]}标记初始化每条序列，确保第一个事件也获得异常分数。逐对相干性分数$\phiMetric$明确\emph{不}包含在输入中——它由$\timegap_i$、$\graphdist_i$和桥接类型确定性确定，包含它会使模型能够“偷懒”，直接依赖这个标量而不学习事件语义模式。

75维原始特征通过学习的线性层投影到$\mathbb{R}^{128}$，然后加上标准正弦位置编码后进入Transformer。

\subsubsection{架构与训练}

我们使用3层Transformer Encoder（d\_model=128，4个注意力头，d\_ff=512，dropout=0.1，约1.5M参数）。在每个自回归步骤中，给定$[e_1, \dots, e_i]$，模型通过多任务预测头预测$e_{i+1}$：

\begin{itemize}
\item \textbf{分类}：边类型（7类CE），源节点类型（3类CE），目标节点类型（3类CE）。
\item \textbf{连续}：源/目标node2higvec（余弦相似度损失），标量节点属性（MSE，归一化到$[0,1]$）。
\item \textbf{时序}：时间间隔桶（10类对数间隔CE：[0-1ms, 1-10ms, $\dots$, 1-24h, >24h]）。
\end{itemize}

总损失为所有分量的等权和。余弦相似度用于节点嵌入，因为稀疏哈希向量上的MSE会驱使模型预测全零；余弦相似度聚焦于语义空间的方向对齐。时间预测使用对数间隔的分类桶，因为事件到达间隔是重尾的——24小时间隔的单个MSE误差产生会破坏训练稳定性的梯度。

训练仅使用从DARPA第2--4天提取的良性链。长度$n < 3$的链被丢弃（单事件异常是TGN的职责）。最大序列长度为64；过长链以种子为中心截断（因果距离离种子最远的事件先丢弃）。训练在一GPU小时内收敛。

\subsection{基于分量eCDF的检测}

推理时，训练好的Transformer计算逐事件预测损失。为将原始损失转换为异常分数，我们使用分量经验累积分布函数（eCDF），在留出的良性验证集上校准：

\begin{enumerate}
\item 对每个损失分量$k \in \{\text{cat}, \text{cont}, \text{time}\}$，在良性验证链上收集逐事件损失分布。
\item 构建$\text{eCDF}_k$：原始损失 $\mapsto$ 百分位数 $\in [0, 1]$。
\item 对于测试事件：异常分数 $= \max(P_{\text{cat}}, P_{\text{cont}}, P_{\text{time}})$。
\item 链异常分数 $= \text{mean}_i (\text{事件异常分数}_i)$。
\end{enumerate}

我们使用分量上的最大值而非总和，因为时间上极端但语义正常的事件（如C2信标在异常间隔出现）可能在总和中被稀释。最大值保留了最敏感的异常维度。这种分离也实现了可解释性：每个告警可被标记为由时间异常、语义异常或结构异常触发。

当链的异常分数超过良性链分数分布的99.9\%分位数时，该链被标记为异常。共享节点且时间接近（$<60$s）的重叠高分链合并为完整的攻击路径。

% §5 度量验证（模型无关）
\label{sec:metric}

本节在独立于任何模型训练的情况下验证因果相干性度量。本节中的所有实验都是模型无关的——不使用Transformer。

\subsection{实验设置}

我们使用DARPA透明计算E3（CADETS）数据集：268,242个节点，29,727,441条边，跨越五天。第2--4天（4月2--4日）作为校准用的良性数据。第6天（4月6日）包含真实攻击（nginx漏洞利用链）。攻击节点的真实标签遵循MAGIC/ThreaTrace标注。参数校准在第2--3天进行；第4天留出用于良性评估。

\subsection{良性链与随机链判别}

\textbf{声明：} $\branchCoherence$能以高判别力分离因果相干的链与随机构造的链。

\textbf{设置。} 我们从第4天（良性）的100个TGN种子中提取链。随机链通过采样随机溯源边作为伪种子并进行无向BFS（不强制因果方向）来构造。对两组链，我们在没有任何模型训练的情况下计算$\branchCoherence$。

<!-- DATA_NEEDED: GAP_S5_BENIGN_RANDOM — Φ_bw分类良性与随机链的AUC，直方图 -->

\textbf{预期。} $\branchCoherence$在分离良性与随机链上达到AUC $> 0.85$。良性链展现出高相干性（真实的系统因果），而随机链展现出低相干性（事件没有因果连接）。这验证了度量在没有任何模型训练的情况下捕获了因果结构。

\subsection{肘部检测验证}

\textbf{声明：} $\chainCoherence$ vs.~$\hoplimit$曲线展现出可检测的肘部，且肘部点产生最优提取。

\textbf{设置。} 我们从100个良性TGN种子中提取链。对每个$\hoplimit \in \{1..10\}$，从相同种子重新提取链并计算均值$\chainCoherence$。我们将$\hoplimit^*$识别为最大负二阶导数点。

<!-- DATA_NEEDED: GAP_S5_ELBOW — Φ vs δ曲线图配肘部标注，δ*处与δ=1和δ=10的链质量对比 -->

\textbf{预期。} 在$\hoplimit^* \approx 2$--$3$处出现清晰肘部。$\hoplimit^*$处提取的链比$\hoplimit=1$处（缺少因果上下文）更完整，比$\hoplimit=10$处（拉入无关事件）噪声更少。

\subsection{攻击链案例研究}

\textbf{声明：} 真实攻击链展现出高$\branchCoherence$，表明尽管攻击在语义上是恶意的，但遵守操作系统的物理规律。

\textbf{设置。} 我们从第6天（攻击日）的TGN种子中提取链。检查包含真实攻击节点的链的$\branchCoherence$。

<!-- DATA_NEEDED: GAP_S5_ATTACK — 攻击链与top良性链的Φ_bw分数对比表 -->

\textbf{预期。} 攻击链具有与良性链可比的$\branchCoherence$。这确认了关注点分离：相干性度量衡量因果完整性（攻击链和良性链都是因果相干的），而Transformer的角色是识别语义异常。

\subsection{与基于规则提取的对比}

\textbf{声明：} 对基于规则提取输出（Holmes）应用$\branchCoherence$得到比我们因果提取的链更低的分数。

\textbf{设置。} 我们从相同的TGN种子应用Holmes式的基于规则的前向/后向追踪，并计算结果路径的$\branchCoherence$。

<!-- DATA_NEEDED: GAP_S5_HOLMES — Φ_bw(我们的) vs Φ_bw(Holmes) vs Φ_bw(随机)对比表 -->

\textbf{预期。} $\branchCoherence(\text{我们的}) > \branchCoherence(\text{Holmes}) > \branchCoherence(\text{随机})$，表明学习的种子选择和尊重因果方向的提取能比基于规则的替代方法产生更紧凑的链。

% §6 端到端评估
\label{sec:evaluation}

本节评估完整的检测流水线，包括自回归Transformer。

\subsection{主要结果：链级APT检测}

\textbf{设置。} 我们在第2--4天良性链上训练自回归Transformer。在第6天（攻击日）测试。所有对比方法使用\emph{相同}的Transformer架构和训练过程——仅链/段提取策略不同。

\textbf{基线。} 我们与代表四种链提取范式的七个基线对比：（1）EagleEye式固定大小时间窗口；（2）Sentient式随机游走场景分割；（3）GET-AID式时序子图划分；（4）SPARSE式基于规则的路径评分；（5）KAIROS窗口级检测；（6）MAGIC实体级检测；（7）SLOT后处理链重建。

\textbf{指标。} 链级精确率、召回率和F1（主要指标）。AUC。调查成本——分析员按降序异常分数检查边以找到所有真实攻击边所需检查的边数量。

<!-- DATA_NEEDED: GAP_S6_MAIN — 所有方法链级F1柱状图，全指标表 (P/R/F1/AUC/InvCost) -->

\textbf{预期。} 我们的因果链提取在链级F1上超过次优方法超过10个百分点的绝对提升。调查成本相比KAIROS窗口级检测减少超过50倍。

\subsection{跨数据集泛化}

\textbf{设置。} 我们将相同流水线应用到DARPA E5（THEIA/Trace）。校准协议在E5良性数据上从头重新运行（协议不变，参数自动适配）。TGN和Transformer在E5良性链上重新训练。

<!-- DATA_NEEDED: GAP_S6_CROSS — 左右并排E3 vs E5肘部曲线，已发现参数表，两者链级F1 -->

\textbf{预期。} 校准协议为E3和E5发现不同的$\hoplimit^*$值（展示了OS专属的适配），但两者上的链级F1始终很高（展示了泛化性）。

\subsection{消融研究}

我们进行六项消融以隔离每个设计组件的贡献：

\begin{enumerate}
\item \textbf{移除因果方向强制：} 将因果追溯替换为无向BFS。预期：最大的F1下降——证明因果方向的重要性。
\item \textbf{移除TGN时序记忆：} 替换为静态边特征进行种子选择。预期：中等下降——证明TGN时序状态的价值。
\item \textbf{移除分支加权相干性：} 仅使用$\chainCoherence$。预期：高分支良性进程（Firefox、编译）导致精确率下降。
\item \textbf{移除Node Profiler：} 仅使用KAIROS原始特征。预期：中等下降。
\item \textbf{移除正弦时间编码：} 替换为标量$\timegap$。预期：轻微下降——但证明多尺度时序的价值。
\item \textbf{将Transformer替换为LSTM和MLP：} 预期：LSTM接近，MLP更差——证明Transformer选择的合理性。
\end{enumerate}

<!-- DATA_NEEDED: GAP_S6_ABLATION — 消融柱状图，逐消融指标表 -->

\subsection{效率分析}

<!-- DATA_NEEDED: GAP_S6_EFFICIENCY — 逐组件延迟分解，内存分析，吞吐量对比 -->

处理一天DARPA E3数据（约500万事件）大约需要10分钟：TGN推理占主导约8分钟，链提取增加约2分钟（CPU），Transformer推理在1秒内完成。这代表了144倍的实时吞吐量乘数。相比KAIROS，我们的方法增加不到20\%的开销。

% 讨论
\section{讨论}
\label{sec:discussion}

\subsection{适用范围与部署模式}

本文方法面向异步带外检测，而非内联阻断。系统在滑动图缓冲区上批量处理审计事件：边级种子识别器先筛出值得调查的事件，链提取和序列模型仅在种子触发时运行。该设计适合持续数小时或数天的 APT 调查场景，牺牲毫秒级阻断能力以换取较高吞吐和可解释的链级输出。

本文评估采用逐日批处理，以便与既有 PIDS 工作公平对比。部署时同一流水线可改为微批流式执行；只要审计事件以时间顺序进入图缓冲区，链提取和 eCDF 校准机制无需改变。

\subsection{局限性}

\textbf{$\branchCoherence$ 不是恶意性分数。}它衡量链是否具有审计支撑的结构一致性，而不是链是否恶意。攻击链和良性链都可能高度一致；恶意性由序列模型的异常分数 $A(\chain)$ 判断。我们通过 \S\ref{sec:metric} 的负控制和敏感性实验验证它确实捕获结构属性，但不把它解释为校准概率。

\textbf{召回受种子识别器约束。}链提取由种子边触发。若攻击步骤完全未被种子识别器浮现，后续链提取无法恢复该步骤。因此端到端召回可分解为种子召回与给定种子覆盖后的链层召回。本文显式报告该分解，以避免把种子识别器能力误归因给链提取器。

\textbf{路径化表示限制。}当前评分对象是以种子为中心的线性化路径，合并后的多个告警链才形成调查图。这适合稀疏、多阶段 APT，但对勒索软件、蠕虫或大规模扫描等分支型攻击不够自然。原生 DAG 评分是直接扩展方向。

\textbf{审计日志覆盖不足的通信通道。}共享内存、部分 \texttt{mmap} IPC 和信号 IPC 在默认审计日志中缺少稳定依赖语义。若攻击主要依赖这些通道，本文方法可能漏检；处理这些通道需要更强的审计 instrumentation 或额外运行时监控。

\textbf{良性训练假设。}本文假设训练数据无攻击污染，这是溯源异常检测中的常见设定。真实部署中可结合鲁棒统计、粗粒度污染过滤或人工确认的良性窗口来放松该假设。

\textbf{跨 engagement 而非完全跨平台。}DARPA E3 与 E5 来自同一计划，虽然场景不同，但审计基础设施相近。因此本文结果只能说明跨 engagement 适配能力；跨 OS、跨审计框架（如 Linux auditd 到 Windows ETW）的泛化仍需进一步验证。

\subsection{崩溃诱导的取证盲区：召回评估的数据上界}
\label{sec:forensic-ceiling}

在核对 CADETS E3 攻击日的原始数据时，我们发现了一个对召回评估有直接影响、且此前工作未曾报告的现象。我们不依赖数据集自带的分片边界，而是对 \texttt{timestampNanos} 做二分查找以精确定位攻击窗口在原始流中的字节偏移。据此我们观测到 04-06 12:08:59--14:01:12 ET 之间存在一段约 1 小时 52 分钟的\emph{物理数据缺口}：该区间内没有任何审计记录。这段缺口与运营事件日志中独立记录的一次故障吻合——``CADETS 目标系统崩溃，原因不明，约 12:20--14:00 ET 停止发布''。

关键在于，该攻击的最后一个活动窗口（约 12:03:50--14:01:32 ET）恰好横跨这段缺口。对照 ground-truth 叙事的末三步——失败的 \texttt{netrecon} 二次探测、\texttt{libdrakon} 载荷投递、以及针对 \texttt{sshd} 的注入尝试——我们在全部约 1330 万条原始记录中做逐字节穷举检索（包括这几步涉及的具体网络地址），命中数为零。也就是说，这三步并非被我们的读取器或种子识别器漏掉，而是\emph{在数据层面根本不存在}。合理的解释是：攻击自身的失稳动作（内存注入尝试）导致主机崩溃，而崩溃连带杀死了审计进程本身，从而摧毁了该攻击后续步骤的溯源证据。无论攻击者是否有意为之，这在客观上构成一种\emph{攻击自毁式的反取证副作用}。

这一现象为召回评估设定了一个与算法无关的数据上界：任何基于溯源图的检测器，无论链提取能力多强，都无法恢复审计日志中不存在的边。因此在该攻击上可达到的召回上界由数据完整性决定（可恢复约 8/11 步），而不仅由种子召回与链层召回决定。据此我们把 \S\ref{sec:evaluation} 的端到端召回分解进一步扩展为
\[
\mathrm{Recall}_{\text{end-to-end}}
= \underbrace{\rho_{\text{data}}}_{\text{数据可恢复性上界}}
\times \mathrm{Recall}_{\text{seed}}
\times \mathrm{Recall}_{\text{chain}\mid\text{seed-covered}},
\]
其中 $\rho_{\text{data}}$ 为落在有效审计覆盖区间内的 ground-truth 步骤占比。我们在报告 CADETS E3 结果时显式标注这三步为``数据层面不可恢复''，仅在可恢复子集上评估链提取质量，以免把数据缺口误算作方法的漏检。

\subsection{未来工作}

未来工作包括四个方向。第一，将当前批处理实现扩展为实时审计日志流处理。第二，设计依赖因果感知的种子识别器，减少对通用 TGN 重构损失的依赖。第三，将 $\branchCoherence$ 从线性路径扩展到原生 DAG 评分，以更好处理分支型攻击。第四，研究跨主机溯源整合；此时依赖边不再完全由单机内核中介，依赖语义也必须相应扩展。

% §8 结论
\label{sec:conclusion}

我们提出了因果相干性度量（$\branchCoherence$）——一个模型无关的、根植于操作系统原理的度量，用于量化从溯源图中提取的事件链的因果结构完整性——以及一个无监督校准协议，仅从良性数据中发现操作系统的自然因果视野。该度量的三分法逐对相干性函数区分身份保持桥接（进程）和有状态桥接（文件和套接字），根植于进程是独占能力而文件是共享资源这一OS原理。校准协议使用良性链上的肘部检测来设置所有参数，无需攻击标签或人工调优，并自动适配不同的OS配置。

在$\branchCoherence$指导下提取的因果链，由一个仅在良性事件序列上训练的自监督自回归Transformer进行验证，在DARPA E3上达到[X]\%的链级F1——超过基于窗口、基于随机游走、基于子图和基于后处理链重建的方法超过10个百分点的绝对F1——同时相比最先进的窗口级检测器减少超过50倍的调查成本。校准协议在E3和E5上发现不同的因果视野，展示了跨平台自动适配能力。

通过将APT检测从“异常在哪里/何时”重新定义为“攻击是哪个因果序列”，我们的方法为安全分析员提供了直接可操作的攻击叙事，而非可疑实体或时间窗口的列表。


% 致谢 (可选)
% \begin{acks}
% 感谢匿名评审者的宝贵意见。
% \end{acks}

\bibliographystyle{ACM-Reference-Format}
\bibliography{references}

\end{document}
